\section{Where It Falls Down}
\label{sec:limits}

\subsection{A small back group means no speed}

Some permutations move the innermost axis. Then every way of breaking the job
up has a swap with nothing behind it. For that swap, $D_{post}=1$.

Now each value we move is a single number on its own. The memory system reads a
whole cache line and uses four bytes of it. There is no way to group the work.

Our second test measures this. We used \hardCases{} permutations picked at
random. In \hardBigCases{} of them a back group survived. We were the best
in-place method on \hardBigWins{} of those \hardBigCases{}. In the other
\hardSmallCases{} no back group survived. We were the worst in-place method on
every one. We took \hardSmallBoxLo{} to \hardSmallBoxHi{}\,ms. The copy took
\hardSmallContLo{} to \hardSmallContHi{}\,ms. We were about
$\hardSmallSlowdown\times$ slower.

We still finish. We still ask for no extra memory. But here the published
in-place methods are several times faster than we are. A real system would send
these cases to one of them. We have not built that. We report the gap instead
of hiding it.

Rank~2 is always in this group. A plain matrix transpose has no back axes to
leave alone.

\subsection{More than one swap loses to the copy}

Section~\ref{sec:fixedpoints} shows a single swap always moves at least
$\modelMinMoved$ of the tensor. So a route with two or more swaps moves at
least as much as a copy does. We are often still the best in-place choice
there. But we are not faster than copying. A better search would not change
that.

\subsection{Planning is slow at high rank}

The search runs over orderings of the axes. It stops as soon as it reaches the
target. So the common cases, which need one swap, come back at once. But a hard
permutation at rank~8 took several seconds to plan.

A route depends only on the shape and the permutation. So we work it out once
and reuse it. A much higher rank would need a search that aims at the target
instead of spreading out in every direction.

\subsection{One GPU}

Every time in this paper comes from one \resDevice{}. The main speedup is
arithmetic. Equation~(\ref{eq:moved}) has no hardware term in it. So it should
hold on other machines. The comparisons against other systems are
measurements, and those need running again elsewhere.

We should also say this. Our rebuild of ITTPD runs its on-chip path about ten
times slower than a plain copy on these shapes. That may be our rebuild rather
than the published method. If so, our lead over it is smaller than we report.
